\section{Discussion}

\paragraph{JEPA representations are adaptable, not immediately discriminative.}
Section~\ref{sec:res_emb} shows a simple pattern.
Frozen ESM-2 embeddings work better before fine-tuning.
JEPA-AMP works better after fine-tuning.
We read this as \emph{representational plasticity}.
JEPA keeps more room for task-specific adaptation, while ESM-2 is stronger
before adaptation.
In practice, JEPA is the better choice when task labels are available for
fine-tuning. ESM-2 may be enough for quick screening without fine-tuning.

\paragraph{AMP-specific pre-training helps, but not by winning every metric.}
JEPA-AMP does not win every benchmark.
On AMPlify, a specialised supervised ensemble remains stronger.
The main value of AMP-domain pre-training is breadth.
The same encoder supports classification, multi-species MIC prediction, and
QMAP regression without changing the backbone.
It is especially strong on high-efficiency E.~coli under homology-aware
evaluation ($+0.098$ PCC over the best prior result).
So the contribution is not ``a better classifier.''
It is a compact AMP-specific representation that works across several tasks,
with clear limits on each task.

\paragraph{Homology-aware benchmarks sharpen biological claims.}
QMAP matters because it reduces the chance that simple near-neighbour
memorisation explains good performance.
Under this stricter protocol, JEPA-AMP is competitive on full E.~coli MIC and
stronger on high-efficiency E.~coli, but the HC50 result clarifies that bacterial
MIC and mammalian haemolysis should not be treated as one interchangeable task.
For AMP discovery, this distinction matters: shared representations may transfer
across bacterial species, but antimicrobial potency and mammalian toxicity require
separate heads and separate supervision.

\paragraph{Controllability is property-specific, not a property of the decoder.}
The generation results draw a clear boundary.
Net charge is directly controllable because it is compositional and receives an
explicit differentiable loss signal; GRAVY and length remain coupled with amino
acid residue choices required to satisfy charge constraints.
MIC conditioning is more indirect: the generator can move JEPA and ESM-2 MIC
scorers toward globally active or inactive regions, but species selectivity fails
because sparse GRAMPA supervision ($\sim$3.5 of 20 bacterial measurements per
peptide) is dominated by broad-spectrum correlations.
So we should not say ``the generator controls peptides'' in a broad sense.
We should say that each requested property needs its own control audit.

\paragraph{Three generation settings on one backbone.}
The comparison in Section~\ref{sec:res_gen_paradigm} gives a clear answer.
The AR decoder is still best for charge control.
NAR works, but it is weaker than AR on both charge and length.
Masked diffusion is the strongest new option: it keeps good charge control and
matches target length exactly in this evaluation.
The v6 model shows a different tradeoff.
Adding more target dimensions helps for helix and pI, but it does not solve
GRAVY, hydrophobic moment, or length control.
In short, decoder choice matters, and more condition variables do not
automatically make every property controllable.

\paragraph{Limitations.}
\begin{enumerate}[noitemsep]
  \item \textbf{GRAVY and length control remain weak.}
        GRAVY correlates with the residue composition required to achieve a
        target charge (K/R are both cationic and hydrophilic); a decoupled
        orthogonal conditioning basis or multi-property optimisation may help.
  \item \textbf{MIC accuracy is moderate.}
        RMSE of 0.62--0.63 in $\log_2$ space corresponds to approximately a
        half-doubling-dilution error; clinical translation requires
        assay-specific validation.
  \item \textbf{MIC-conditioned generation is oracle-mediated.}
        The ESM-2 scorer reduces single-oracle circularity, but both are
        in-silico validators. Wet-lab synthesis and antimicrobial assays are
        required before any generated sequence can be described as
        biologically active.
  \item \textbf{Species-selective control is a data limitation.}
        GRAMPA's sparse per-sequence bacterial coverage makes fine-grained
        selective activity a supervision gap, not an architecture failure.
  \item \textbf{Single-split MIC evaluation.}
        The GRAMPA formal evaluation uses one seed-42 split; multi-seed
        or cross-validation estimates are not yet reported.
  \item \textbf{Positional embedding length limit.}
        Sequences longer than 50 AAs are truncated; future encoders should
        remove this constraint.
  \item \textbf{Corpus annotation bias.}
        The 868k pre-training corpus relies on automated database curation;
        a community benchmark beyond AMPlify would enable definitive comparisons.
\end{enumerate}

\paragraph{Path to wet-lab extension.}
The next wet-lab step is simple.
We should synthesise a small and diverse set of generated peptides plus matched
scrambled controls, then test MIC on at least E.~coli and S.~aureus together
with a haemolysis or cytotoxicity screen.
This would test ranking quality, global activity shifts, and safety filtering
without making stronger claims than this paper can support.
We have already generated 20 top-ranked candidates for this future screen.

\section{Conclusion}

JEPA-AMP shows that AMP-specific self-supervised learning is useful mainly as a
reusable backbone for dry-lab AMP prioritisation.
The same encoder approaches curated AMP classifiers, improves MIC prediction
over ESM-2, and remains competitive under homology-aware QMAP evaluation.
The plasticity analysis also shows when JEPA is preferable:
it has more fine-tuning headroom than ESM-2, while ESM-2 frozen embeddings
may still be enough for rapid screening.
The generation study draws a complementary boundary: charge control is reliable
via the dual-pathway decoder, while GRAVY, length, and species-selective MIC
control remain limited.
Comparing AR, NAR, masked diffusion, and extended conditioning on the same
backbone shows that some limits come from the decoder, while others come from
the target properties themselves.
Together, these results support a cautious computational pipeline for ranking
AMP candidates.
Model predictions should be treated as prioritisation evidence, not as
biological validation.

\bibliographystyle{plainnat}
\bibliography{refs}
